\section{Corrigés}
\mxneed{6}
\begin{solution}[Exercice \ref{mada-t11-s-01:img}]
(1) $f(0)=1$ ; $f(-2)=4+8+1=13$ ; $f(\frac12)=\frac14-2+1=-\frac34$. (2) $x^2-4x+1=-2\iff x^2-4x+3=0$, de racines $1$ et $3$ : les antécédents de $-2$ sont $1$ et $3$. (3) $x^2-4x+5=0$ a pour discriminant $16-20=-4<0$ : $-4$ n'a aucun antécédent.
\end{solution}
\begin{solution}[Exercice \ref{mada-t11-s-01:quot}]
(1) $g(0)=-\frac13$ ; $g(4)=9$ ; $g(-1)=\frac{-1}{-4}=\frac14$. (2) $g(x)=2$ avec $x\ne3$ donne $2x+1=2x-6$, impossible : $2$ n'a pas d'antécédent (on a $g(x)=2+\frac{7}{x-3}$, jamais égal à $2$). $g(x)=1$ donne $2x+1=x-3$, soit $x=-4$, valeur admise ($\ne3$).
\end{solution}
\begin{solution}[Exercice \ref{mada-t11-s-01:dom1}]
(1) $\R$. (2) $\R\setminus\ens{7}$. (3) $\R\setminus\ens{-4;4}$. (4) $3x+6\ge0$ : $\intfo{-2}{+\infty}$. (5) $8-2x\ge0$ : $\intof{-\infty}{4}$. (6) $x+3>0$ : $\intoo{-3}{+\infty}$.
\end{solution}
\begin{solution}[Exercice \ref{mada-t11-s-01:dom2}]
$x^2+4>0$ pour tout $x$ : $D_f=\R$. $x^2-4=0$ en $\pm2$ : $D_g=\R\setminus\ens{-2;2}$. Sur $D_g$, $g(x)=\frac{x+2}{(x-2)(x+2)}=\frac{1}{x-2}$. L'expression $\frac{1}{x-2}$ est définie en $-2$, pas $g$ : les deux ensembles de définition diffèrent, d'où la règle de les chercher sur la formule d'origine.
\end{solution}
\begin{solution}[Exercice \ref{mada-t11-s-01:var1}]
$x\mapsto-2x+5$ : décroissante sur $\R$ (pente négative). $3x^2$ : croissante sur $\intfo{0}{+\infty}$, décroissante sur $\intof{-\infty}{0}$. $-x^2$ : décroissante sur $\intfo{0}{+\infty}$ (multiplication du carré par $-1$).
\end{solution}
\begin{solution}[Exercice \ref{mada-t11-s-01:par1}]
(1) Paire : $f(-x)=f(x)$. (2) Impaire : $(-x)^3+5(-x)=-(x^3+5x)$. (3) Ni l'une ni l'autre : $f(1)=4$, $f(-1)=0$. (4) $D=\R\setminus\ens{-1;1}$, symétrique, $f(-x)=\frac{-x}{x^2-1}=-f(x)$ : impaire. (5) Paire : $\abs{-x}=\abs x$ et $(-x)^2=x^2$. (6) $D=\intfo{0}{+\infty}$ non symétrique : ni l'une ni l'autre. (7) $1\notin D$ mais $-1\in D$ : ni l'une ni l'autre.
\end{solution}
\begin{solution}[Exercice \ref{mada-t11-s-01:per1}]
$31=30+1$ : $u(31)=u(1)=3$. $-4+6=2$ : $u(-4)=u(2)=2$. $300=3\times100$ : $u(300)=u(0)=1$.
\end{solution}
\begin{solution}[Exercice \ref{mada-t11-s-01:per2}]
$\frac{2\pi}{3}$ pour $f$ ; $\frac{2\pi}{1/2}=4\pi$ pour $g$ ; $\frac{2\pi}{2}=\pi$ pour $h$ (la phase $\frac\pi3$ ne change pas la période).
\end{solution}
\begin{solution}[Exercice \ref{mada-t11-s-01:dom3}]
(1) $x^2-6x+8=(x-2)(x-4)\ge0$ : $\intof{-\infty}{2}\cup\intfo{4}{+\infty}$. (2) $5-x\ge0$ et $x\ne\pm3$ : $\intoo{-\infty}{-3}\cup\intoo{-3}{3}\cup\intof{3}{5}$. (3) $\abs x\ne3$ : $\R\setminus\ens{-3;3}$. (4) Il faut $x\ne-4$ et $\frac{x-1}{x+4}\ge0$ : positif pour $x<-4$ et pour $x\ge1$, négatif entre $-4$ et $1$ : $\intoo{-\infty}{-4}\cup\intfo{1}{+\infty}$.
\end{solution}
\begin{solution}[Exercice \ref{mada-t11-s-01:var2}]
Soit $a<b<0$. Alors $\frac1a-\frac1b=\frac{b-a}{ab}$ avec $b-a>0$ et $ab>0$ (deux réels négatifs), donc $\frac1a>\frac1b$ : l'ordre est renversé, la fonction est strictement décroissante sur $\intoo{-\infty}{0}$.
\end{solution}
\begin{solution}[Exercice \ref{mada-t11-s-01:var3}]
Comme $-1\le2\le6$ et que $f$ est croissante, $3\le f(2)\le9$. On ne peut pas donner la valeur de $f(2)$, ni dire si $f$ est strictement croissante, ni si sa courbe est continue.
\end{solution}
\begin{solution}[Exercice \ref{mada-t11-s-01:var4}]
(1) Pour $a\le b$ : $f(a)\ge f(b)$ et $g(a)\ge g(b)$, d'où $f(a)+g(a)\ge f(b)+g(b)$. (2) $fg(x)=x^2$ n'est pas croissante sur $\intoo{-\infty}{0}$ : $-3<-1$ mais $9>1$. On ne multiplie pas des inégalités sans précaution.
\end{solution}
\begin{solution}[Exercice \ref{mada-t11-s-01:par2}]
Le domaine de $f$ est $\intfo{-2}{+\infty}$ : $3$ y est, pas $-3$. Le domaine de $g$ est $\R\setminus\ens{5}$ : $-5$ y est, pas $5$. Aucun n'est symétrique par rapport à $0$.
\end{solution}
\begin{solution}[Exercice \ref{mada-t11-s-01:axe}]
$f(3+h)=(3+h)^2-6(3+h)+5=h^2-4$ ne change pas quand $h\to-h$ : $f(6-x)=f(x)$. Pour $g$, $a=-1$ : $g(-2-x)=\abs{-1-x}=\abs{x+1}=g(x)$.
\end{solution}
\begin{solution}[Exercice \ref{mada-t11-s-01:centre}]
$D=\R\setminus\ens{1}$ et $2-x\in D$ si $x\in D$. $f(2-x)=2+\frac{3}{1-x}=2-\frac{3}{x-1}$ ; avec $f(x)=2+\frac{3}{x-1}$, la somme vaut $4=2\times2$ : $\Omega\pt{1;2}$ est un centre de symétrie.
\end{solution}
\begin{solution}[Exercice \ref{mada-t11-s-01:par3}]
$(fg)(-x)=f(x)\bigl(-g(x)\bigr)=-(fg)(x)$ : impaire. $f(-x)^2=f(x)^2$ : $f^2$ paire. Pour $f+g$, avec $f=1$ et $g(x)=x$ : $(f+g)(1)=2$ et $(f+g)(-1)=0$ : ni paire ni impaire. Si $f$ est paire et impaire, $f(x)=f(-x)=-f(x)$, donc $2f(x)=0$ pour tout $x$.
\end{solution}
\begin{solution}[Exercice \ref{mada-t11-s-01:per3}]
$f(12)=f(2)=4$ ; $f(-1)=f(4)=16$ ; $f(\num{17,5})=f(\num{2,5})=\num{6,25}$.
\end{solution}
\begin{solution}[Exercice \ref{mada-t11-s-01:per4}]
(1) $\cos^2(x+\pi)=(-\cos x)^2=\cos^2x$ ; et $\cos^2(-x)=\cos^2x$. (2) $\sin$ et $\cos$ ont la période $2\pi$ donc leur somme aussi. (3) Les deux termes ont la période $2\pi$ et $f$ est paire : on étudie $f$ sur $\intff{0}{\pi}$.
\end{solution}
\begin{solution}[Exercice \ref{mada-t11-s-01:per5}]
$f(x-T)=f\bigl((x-T)+T\bigr)=f(x)$. De $f(x+T)=f(x)$ on tire de proche en proche $f(x+3T)=f(x+2T)=f(x+T)=f(x)$. Enfin $f(x+S+T)=f\bigl((x+S)+T\bigr)=f(x+S)=f(x)$.
\end{solution}
\begin{solution}[Exercice \ref{mada-t11-s-01:imp}]
$f(0)=-f(-0)=-f(0)$ donc $f(0)=0$. Soit $a<b\le0$ : $0\le-b<-a$, la croissance sur $\intfo{0}{+\infty}$ donne $f(-b)\le f(-a)$, soit $-f(b)\le-f(a)$, donc $f(a)\le f(b)$.
\end{solution}
\begin{solution}[Exercice \ref{mada-t11-s-01:axper}]
Pour $h$ réel : $f(\frac T2+h)=f(h-\frac T2)$ (période $T$) $=f(\frac T2-h)$ (parité). Avec $x=\frac T2+h$ on a $f(T-x)=f(x)$, c'est-à-dire $f(2a-x)=f(x)$ pour $a=\frac T2$.
\end{solution}
\begin{solution}[Hors programme : minimalité]
$\cos^2x=\frac12+\frac12\cos 2x$ : la plus petite période de $\cos 2x$ est $\frac{2\pi}{2}=\pi$ ; ajouter la constante $\frac12$ et diviser par $2$ ne la change pas.
\end{solution}
\begin{solution}[Exercice \ref{mada-t11-s-01:recon}]
(1) $f(-\num{0,5})=f(\num{0,5})=1$ ; $f(3)=f(-1)=f(1)=2$ ; $f(10)=f(2)=2$ ; $f(\num{7,25})=f(\num{3,25})=f(-\num{0,75})=f(\num{0,75})=\num{1,5}$. (2) Sur $\intff{-1}{1}$, $f(x)=2\abs x$ ; sur $\intff{1}{3}$, $f(x)=2$ ; la courbe se répète tous les $4$. (3) Le maximum est $2$, atteint sur chaque intervalle $\intff{1+4k}{3+4k}$, $k$ entier.
\end{solution}
\begin{solution}[Exercice \ref{mada-t11-s-01:roue}]
(1) $\frac{2\pi}{\pi/10}=20$ : un tour complet dure $20$ minutes. (2) $h(0)=1$ ; $h(10)=11$ ; $h(35)=h(15)=6-5\cos\frac{3\pi}{2}=6$ ; $h(50)=h(10)=11$. (3) $h(10+u)=6-5\cos(\pi+\frac{\pi u}{10})=6+5\cos\frac{\pi u}{10}$ et $h(10-u)=6-5\cos(\pi-\frac{\pi u}{10})=6+5\cos\frac{\pi u}{10}$ : la courbe est symétrique par rapport à $t=10$, le sommet. (4) $\cos\frac{\pi t}{10}=0$ donne $t=5$ ou $t=15$.
\end{solution}
